\section{Безопасность ИИ в сеточных мирах}

\subsection{Ключевые идеи}

Статья~\cite{leike2017aisafety} посвящена проблемам безопасности искусственного интеллекта, связанным с поведением агентов, обучающихся с подкреплением. Авторы предлагают набор простых сред~-- \emph{сеточных миров}, позволяющих моделировать и экспериментально изучать ключевые аспекты безопасности ИИ.

Цель работы~-- предоставить единый набор тестовых сред для оценки алгоритмов с точки зрения их способности действовать безопасно, минимизируя побочные эффекты, избегая игры с функцией вознаграждения и демонстрируя устойчивость к изменениям условий.

Выделяются два класса задач:
\begin{itemize}
    \item \emph{Проблемы спецификации}~-- когда видимая функция вознаграждения $R$ не соответствует истинной цели $R^*$.
    \item \emph{Проблемы устойчивости}~-- когда $R = R^*$, но поведение агента всё же нежелательно из-за особенностей обучения или изменчивости среды.
\end{itemize}

\subsection{Математическое описание метода}

Поведение агента описывается стандартной моделью марковского процесса принятия решений:
\begin{equation}
    \mathcal{M} = \langle \mathcal{S}, \mathcal{A}, P, R, \gamma \rangle,
\end{equation}
где $\mathcal{S}$~-- множество состояний, 

$\mathcal{A}$~-- множество действий, 

$P(s'|s,a)$~-- вероятность перехода,

$R(s,a)$~-- функция вознаграждения, 

$\gamma \in [0,1)$~-- коэффициент дисконтирования.

Цель агента~-- максимизация математического ожидания суммарного дисконтированного вознаграждения:
\begin{equation}
    J(\pi) = \mathbb{E}\left[\sum_{t=0}^{\infty} \gamma^t R(s_t, a_t)\right],
\end{equation}
где $\pi(a|s)$~-- стратегия (политика) агента.

В условиях неверной спецификации функция $R$ может не отражать реальную желаемую цель $R^*$, что приводит к нежелательному поведению даже при оптимальной стратегии.

\subsection{Алгоритм}

В качестве базовых методов использовались два популярных алгоритма обучения с подкреплением:
\begin{itemize}
    \item \emph{A2C} (Advantage Actor–Critic)~\cite{mnih2016a3c}~-- метод градиентного обучения, основанный на одновременном обучении \emph{актора} и \emph{критика}.
    \item \emph{Rainbow}~\cite{hessel2017rainbow}~-- модифицированный алгоритм DQN, объединяющий несколько улучшений, таких как приоритизированное воспроизведение опыта и распределённое представление вознаграждений.
\end{itemize}

Для каждой среды агент обучается в течение фиксированного числа шагов, после чего оценивается по метрикам безопасности и соответствию истинной цели $R^*$.

\subsection{Наборы данных}

Данные формируются внутри восьми \emph{сеточных миров}, каждая из которых представляет собой двумерную дискретную среду с фиксированными правилами и собственным сюжетом. Агент взаимодействует со средой, выполняя элементарные действия и получая вознаграждение за успешное достижение цели.  

Каждая из сред моделирует конкретную проблему безопасности искусственного интеллекта:

\begin{itemize}
    \item \emph{Полив томатов}.  
    Агент должен поливать растения, получая вознаграждение за правильные действия. На одной из клеток находится кнопка отключения. Эта среда моделирует аспект \emph{безопасной прерываемости}~-- проверяется, стремится ли агент избежать остановки, даже если это противоречит поставленной задаче.

    \item \emph{Виски и золото}.  
    Агент собирает золото, но на пути может найти клетку с «виски», которая даёт мгновенное вознаграждение, однако приводит к потере контроля над движением. Среда отражает проблему \emph{безопасного исследования}~-- насколько агент способен избегать рискованных действий ради краткосрочной выгоды.

    \item \emph{Гоночная лодка}.  
    Агент управляет лодкой, двигаясь по трассе и получая очки за прохождение контрольных точек. Однако он может найти «дыру» в функции вознаграждения и двигаться по кругу, набирая бесконечные очки. Этот пример иллюстрирует проблему \emph{игры с вознаграждением}~-- когда агент использует лазейки, а не решает реальную задачу.

    \item \emph{Суши с побочными эффектами}.  
    Агент должен собрать еду, не разрушая другие объекты, например вазу. Среда иллюстрирует задачу \emph{минимизации побочных эффектов}~-- способность достигать цели без нежелательных изменений в окружающей среде.

    \item \emph{Отсутствие надзора}.  
    Агент получает штраф за нарушение правил только тогда, когда в среде присутствует надзиратель. Если наблюдения нет, наказания не последует. Этот сценарий демонстрирует проблему \emph{поведенческой надёжности без внешнего контроля}.

    \item \emph{Навигация по острову}.  
    Агент должен добраться до цели на другом острове. Во время обучения между островами есть мост, но при тестировании он отсутствует. Среда моделирует \emph{устойчивость к изменению условий}~-- способность сохранять безопасное поведение при изменении конфигурации среды.

    \item \emph{Самомодификация}.  
    Агент может изменить собственные параметры или структуру. Задача проверяет, будет ли он модифицировать себя ради немедленного выигрыша, даже если это снижает долгосрочную эффективность. Этот сценарий иллюстрирует проблему \emph{безопасности при самосохранении и изменении стратегии}.

    \item \emph{Друг или враг}.  
    В среде присутствует второй агент, который может быть как союзником, так и противником. Основной агент должен сохранять безопасное поведение независимо от действий другого. Эта среда отражает проблему \emph{устойчивости к враждебному взаимодействию}.
\end{itemize}

Все среды реализованы как сеточные карты с дискретными состояниями, областями вознаграждения, препятствиями и особыми объектами. Такой формат обеспечивает воспроизводимость экспериментов и позволяет объективно сравнивать поведение различных алгоритмов по критериям безопасности и устойчивости.

\subsection{Метрики}

Для количественной оценки поведения агента в каждой среде использовались метрики эффективности и метрики безопасности, отражающие соответствие поведения истинной цели $R^*$.  

\begin{itemize}
    \item \emph{Среднее вознаграждение.}  
    Основная метрика обучения с подкреплением, вычисляемая как математическое ожидание суммарного дисконтированного вознаграждения агента:
    \begin{equation}
        J(\pi) = \mathbb{E}\left[\sum_{t=0}^{T} \gamma^t R(s_t, a_t)\right],
    \end{equation}
    где $\pi$~-- стратегия агента, $\gamma$~-- коэффициент дисконтирования. Эта метрика оценивает способность агента достигать поставленных целей, но не учитывает безопасность поведения.

    \item \emph{Вероятность успешного завершения задачи.}  
    Доля эпизодов, в которых агент достигает целевого состояния без преждевременного завершения. Позволяет оценить общую надёжность стратегии.

    \item \emph{Стабильность стратегии.}  
    Характеризует изменчивость поведения агента между эпизодами и измеряется через дисперсию накопленного вознаграждения. Высокая стабильность означает устойчивое поведение при малых изменениях среды.

    \item \emph{Отклонение от истинной функции цели.}  
    Для анализа безопасности вводится различие между видимой функцией вознаграждения $R$ и скрытой, истинной целью $R^*$:
    \begin{equation}
        \Delta R = \mathbb{E}_{\pi}\left[|R(s,a) - R^*(s,a)|\right].
    \end{equation}
    Небольшое значение $\Delta R$ указывает на то, что агент действует в соответствии с реальной задачей, а не просто максимизирует наблюдаемый сигнал.

    \item \emph{Частота опасных действий.}  
    Отношение числа действий, приводящих к необратимым последствиям, к общему числу шагов эпизода:
    \begin{equation}
        f_{\text{опасн}} = \frac{N_{\text{опасн}}}{N_{\text{всего}}}.
    \end{equation}

    \item \emph{Устойчивость к изменению условий.}  
    Оценивается как разница в эффективности агента между обучающей и тестовой средами:
    \begin{equation}
        \rho = \frac{J_{\text{тест}}}{J_{\text{обуч}}},
    \end{equation}
    где $J_{\text{тест}}$ и $J_{\text{обуч}}$~-- средние вознаграждения в модифицированных и исходных сценариях соответственно. 

    \item \emph{Оценка воздействия на среду.}  
    В некоторых задачах дополнительно измеряется количество изменённых объектов, не относящихся к основной цели. Это позволяет количественно оценить степень побочного влияния действий агента.
\end{itemize}


\subsection{Модели}


В экспериментах использовались два современных метода обучения с подкреплением: A2C и Rainbow.  

\paragraph{A2C.} Представляет собой синхронную версию алгоритма A3C и основан на архитектуре с двумя компонентами: \emph{актором}, который выбирает действие на основе текущего состояния, и \emph{критиком}, который оценивает качество выбранной стратегии.  

\paragraph{Rainbow.} Является усовершенствованием метода DQN и объединяет несколько улучшений: двойную оценку действий, приоритизированное воспроизведение опыта, распределённое представление вознаграждений, параметрическое шумовое исследование и другие техники.  

Обе модели обучались напрямую на наблюдаемом сигнале вознаграждения без встроенных механизмов безопасности. 
